
\section{Related Work}

The classical literature established the difficulty of ocular contamination and motivated automated statistical detection directly from EEG. Early work emphasized artifact removal and statistical screening rather than blink-region benchmarking per se, but it already showed that extreme-value behavior, channel context, and morphology matter \citep{croft2000removal,klein2013reliable}. This historical backdrop explains why simple one-value global thresholds can be attractive, yet brittle.

A first important stream comprises single-channel or few-channel threshold-centered detectors. \citet{chang2016detection} detect blink artifacts from a single prefrontal channel using digital filtering and rule-based decisions. \citet{tran2021detection} present an explicitly thresholding-based EEG blink detector and compute a minimum amplitude threshold from transformed signal magnitude and signal standard deviation. \citet{zhang2023method} move toward real-time frontal EEG detection with short windows and decision logic around potential blink boundaries. \citet{wang2022multidimensional} extend thresholding to harder clinical settings by optimizing multidimensional features in the presence of epileptiform discharges, and later work based on higher-order cumulants continues the move toward short-segment single-channel detection \citep{wang2025sliding}.

A second stream contains hybrid thresholding pipelines. BLINKER first identifies candidate intervals with a mean-based standard-deviation threshold and then rejects implausible candidates using robust amplitude-distribution criteria derived from the median absolute deviation \citep{kleifges2017blinker}. \citet{cao2021unsupervised} use cascaded thresholding over amplitude, amplitude displacement, and cross-channel correlation before applying a Gaussian mixture model. \citet{agarwal2019blink} self-learn user-specific profiles in an unsupervised manner, while \citet{valderrama2018automatic} use iterative template matching and automatic threshold/template estimation for suppression in single-channel recordings. \citet{guttmann2019new} further emphasize adaptation to inter- and intra-subject variability, underscoring that fixed thresholds are often not enough.

A third stream is the recent learning-based and multimodal literature. The recent review of deep learning in blink detection shows that the field is increasingly data-driven, yet also constrained by data imbalance, device heterogeneity, and generalization challenges \citep{xiong2025review}. For a threshold paper, this trend matters for two reasons. First, it confirms that threshold methods remain useful as transparent baselines. Second, it motivates a bridge strategy in which LLMs or other generative tools are not used to replace evaluation, but to propose candidate formulas that remain interpretable and falsifiable \citep{gottweis2025toward}.

What is still missing is a benchmark that compares threshold families themselves under shared data and shared scoring. Prior papers generally optimize their own end-to-end pipeline and report strong performance in their own settings, but they do not isolate threshold choice as the controlled independent variable. That is the gap targeted here.
